\documentclass[11pt,letterpaper]{article}
\usepackage[margin=1in]{geometry}
\usepackage[T1]{fontenc}
\usepackage{amsmath,amssymb,microtype}
\usepackage[hidelinks]{hyperref}
\pagestyle{empty}
\setlength{\parindent}{0pt}
\setlength{\parskip}{0.7em}
\hypersetup{pdftitle={Cover letter to Mathematics of Computation},pdfauthor={Haotian Zhong}}
\begin{document}
\textbf{Haotian Zhong}\\
Department of Mathematics, School of Sciences\\
Shanghai University, Shanghai 200444, P. R. China\\
\href{mailto:zhonghaotian1219@gmail.com}{zhonghaotian1219@gmail.com}\\
ORCID: 0009-0009-5567-1315

September 29, 2026

Editors of \textit{Mathematics of Computation}\\
American Mathematical Society

Dear Editors,

I submit \textbf{``Frequency Saturation and Certified Computation of Relaxation Functionals''} for consideration in \textit{Mathematics of Computation}. The paper determines the frequency information required to recover linear functionals of positive relaxation measures and gives a finite, verified numerical method for computing the complete range allowed by noisy observations.

For resistance--capacitance responses on a compact interval of positive relaxation times, the exact threshold for uniform noiseless identification of references with at most $K$ atoms is $K+2$ distinct frequencies when resistance and inductance are unknown, and $K+1$ when they are known. At the calibrated threshold, a fixed finite design attains the same sharp worst-class functional accuracy order as the entire positive-frequency response. Positive moment-matched spectra yield all-frequency lower bounds, including against target-aware adaptive queries, and finite-design polynomial reproduction gives matching upper bounds. The rate resolves local target regularity and cluster multiplicity uniformly through colliding atoms and vanishing masses.

The computational contribution is an exact-arithmetic refinement algorithm for both endpoints of the original continuous feasible range. On strictly feasible rational observation boxes, it terminates at a prescribed numerical excess. Verified finite optimization gaps, target approximation, and continuum correction enter separate error budgets; feasible witnesses also provide an absolute--relative stopping criterion. The paper includes an explicit algorithm, its complete termination proof, and reproducible continuous-domain certificates.

The same functional modulus gives matching expected acquisition costs for prescribed terminal width and confidence under a Gaussian variance--effort law. The lower bound permits adaptive frequency, precision, and stopping decisions; a class-tuned expanding design invoking the certified range calculation attains the order and completes pointwise over all finite positive spectra. This joins analysis of a restricted information class to an executable, accuracy-controlled inverse computation, addressing approximation, information complexity, and verified numerical methods.

The manuscript is original, unpublished, and not under consideration elsewhere. I declare no conflicts of interest. The accompanying archive provides exact inputs, verification code, certificates, and reproduction instructions. I suggest Professor Markus Bachmayr as a handling editor, given his expertise in approximation, adaptive methods, and inverse problems.

Sincerely,\\[0.35em]
Haotian Zhong
\end{document}
